% !TEX root = ../main.tex
%% Corrigés — Chapitre 17 — Probabilités sur un univers fini
%% Résultats vérifiés par verifications/chapitre_17.py.

\corrigesChapitre{17}{Probabilités sur un univers fini}
\corrigesSouples

\begin{blocLivre}{cor-17-1}
\begin{corrige}{1}{Modéliser une expérience}
\demarche{Univers équiprobable : compter les cas favorables (chapitre 16).}
\begin{questions}
  \item $\Omega=\ens{1,\dots,6}^{2}$, $36$ issues équiprobables. La somme
    vaut $7$ pour les $6$ couples $(1,6),\dots,(6,1)$ :
    $P=\dfrac{6}{36}=\res{\dfrac16}$.
  \item $\Omega=\ens{P,F}^{3}$, $8$ issues ; exactement deux piles :
    $\dbinom32=3$ issues, $P=\res{\dfrac38}$.
\end{questions}
\end{corrige}
\end{blocLivre}

\begin{blocLivre}{cor-17-2}
\begin{corrige}{2}{Les règles de calcul}
\demarche{Union, différence, contraire : trois formules, aucun piège.}
$P(A\cup B)=\dfrac12+\dfrac25-\dfrac15=\res{\dfrac{7}{10}}$ ;
$P(A\setminus B)=P(A)-P(A\cap B)=\res{\dfrac{3}{10}}$ ;
$P\bigl(\overline A\cap\overline B\bigr)
=P\bigl(\overline{A\cup B}\bigr)=1-\dfrac{7}{10}=\res{\dfrac{3}{10}}$.
\end{corrige}
\end{blocLivre}

\begin{blocLivre}{cor-17-3}
\begin{corrige}{3}{Tirages sans remise}
\demarche{Multiplier le long de l'arbre : composées.}
\begin{questions}
  \item $P(R_1\cap R_2)=P(R_1)\,P(R_2\mid R_1)
    =\dfrac35\times\dfrac24=\res{\dfrac{3}{10}}$.
  \item Après une rouge, il reste $2$ rouges sur $4$ boules :
    $P(R_2\mid R_1)=\res{\dfrac12}$.
\end{questions}
\end{corrige}
\end{blocLivre}

\begin{blocLivre}{cor-17-4}
\begin{corrige}{4}{La formule des probabilités totales}
\demarche{Les urnes forment un système complet : on somme les feuilles
rouges de l'arbre.}
Avec le système complet $\ens{U_A,U_B}$ :
\[
P(R)=P(R\mid U_A)P(U_A)+P(R\mid U_B)P(U_B)
=\frac23\cdot\frac12+\frac13\cdot\frac12=\res{\frac12} .
\]
L'arbre a deux étages : choix de l'urne ($\frac12$, $\frac12$), puis
couleur ($\frac23$ ou $\frac13$ vers le rouge).
\end{corrige}
\end{blocLivre}

\begin{blocLivre}{cor-17-5}
\begin{corrige}{5}{La formule de Bayes}
\demarche{Totales pour le dénominateur, Bayes pour remonter du test à la
maladie.}
\begin{questions}
  \item $P(+)=0{,}01\times0{,}95+0{,}99\times0{,}10
    =0{,}0095+0{,}099=\res{0{,}1085}$.
  \item $P(M\mid+)=\dfrac{0{,}0095}{0{,}1085}=\dfrac{19}{217}
    \approx\res{8{,}8\,\%}$ : moins d'un positif sur dix est réellement
    malade ! La maladie est si rare que les faux positifs ($10\,\%$ de
    $99\,\%$) écrasent les vrais positifs ($95\,\%$ de $1\,\%$).
\end{questions}
\end{corrige}
\end{blocLivre}

\begin{blocLivre}{cor-17-6}
\begin{corrige}{6}{Uniforme et Bernoulli}
\demarche{König-Huygens fait tous les calculs de variance.}
\begin{questions}
  \item $\mathrm{E}(X)=\dfrac{1+\dots+6}{6}=\res{\dfrac72}$ ;
    $\mathrm{E}(X^{2})=\dfrac{91}{6}$, d'où
    $\mathrm{V}(X)=\dfrac{91}{6}-\dfrac{49}{4}=\res{\dfrac{35}{12}}$.
  \item $\mathrm{E}(Y)=1\cdot p+0=p$ et
    $\mathrm{E}(Y^{2})=p$ (car $Y^{2}=Y$) :
    $\mathrm{V}(Y)=p-p^{2}=\res{p(1-p)}$, maximale en
    $p=\dfrac12$ (valeur $\frac14$) : l'incertitude est maximale pour une
    pièce équilibrée.
\end{questions}
\end{corrige}
\end{blocLivre}

\begin{blocLivre}{cor-17-7}
\begin{corrige}{7}{La loi binomiale}
\demarche{k succès parmi n répétitions : combinaison fois probabilité d'une
issue.}
\begin{questions}
  \item $P(X=2)=\dbinom52\Bigl(\dfrac13\Bigr)^{2}\Bigl(\dfrac23\Bigr)^{3}
    =\res{\dfrac{80}{243}}$ ; la somme des $P(X=k)$ vaut
    $\bigl(\frac13+\frac23\bigr)^{5}=1$ — c'est la formule du binôme
    (chapitre 16).
  \item $\mathrm{E}(X)=np=\res{\dfrac53}$,
    $\mathrm{V}(X)=np(1-p)=\res{\dfrac{10}{9}}$.
  \item Le nombre de succès lors de $5$ répétitions indépendantes d'une
    expérience réussissant avec probabilité $\frac13$.
\end{questions}
\end{corrige}
\end{blocLivre}

\begin{blocLivre}{cor-17-8}
\begin{corrige}{8}{Espérance : transfert et linéarité}
\demarche{Le transfert évite de chercher la loi de f(X).}
\begin{questions}
  \item $\mathrm{E}\bigl(X^{2}\bigr)=\dfrac{1+4+\dots+36}{6}
    =\res{\dfrac{91}{6}}\neq\Bigl(\dfrac72\Bigr)^{2}$ : l'écart est
    précisément la variance.
  \item $\mathrm{E}(2X+3)=2\times\dfrac72+3=\res{10}$.
  \item $\ind A$ vaut $1$ avec probabilité $P(A)$ et $0$ sinon :
    $\mathrm{E}(\ind A)=P(A)$ — les probabilités sont des espérances
    d'indicatrices, le pont entre les chapitres.
\end{questions}
\end{corrige}
\end{blocLivre}

\begin{blocLivre}{cor-17-9}
\begin{corrige}{9}{Indépendance : exemples et pièges}
\demarche{Vérifier la formule produit — pour toutes les sous-familles.}
\begin{questions}
  \item $P(A)=\dfrac12$, $P(B)=\dfrac12$ ; et $A\cap B$ (premier pair
    \emph{et} somme paire, i.e. les deux pairs) : $\dfrac{9}{36}=\dfrac14
    =P(A)P(B)$ : \res{\text{indépendants}}.
  \item $P(A)=P(B)=P(C)=\dfrac12$ et chaque intersection double vaut
    $\dfrac14$ : deux à deux indépendants. Mais
    $A\cap B\subseteq C$, donc
    $P(A\cap B\cap C)=\dfrac14\neq\dfrac18=P(A)P(B)P(C)$ : la famille
    n'est \res{\text{pas mutuellement indépendante}}.
\end{questions}
\end{corrige}
\end{blocLivre}

\begin{blocLivre}{cor-17-10}
\begin{corrige}{10}{Loi conjointe et lois marginales}
\demarche{Les marginales se lisent en sommant lignes et colonnes — et ne
suffisent pas à reconstruire la table.}
\begin{questions}
  \item Chaque marginale vaut $\bigl(\frac12,\frac12\bigr)$ : lois de
    Bernoulli $\mathcal{B}\bigl(\frac12\bigr)$ ; et
    $\frac14=\frac12\times\frac12$ dans chaque case :
    \res{\text{indépendantes}}.
  \item Mêmes marginales $\bigl(\frac12,\frac12\bigr)$, mais
    $P(X=0,Y=1)=0\neq\frac14$ : \res{\text{dépendantes}} (ici $Y=X$ !).
    Deux couples peuvent avoir les mêmes marginales et des lois conjointes
    radicalement différentes.
\end{questions}
\end{corrige}
\end{blocLivre}

\begin{blocLivre}{cor-17-11}
\begin{corrige}{11}{Minimum et maximum de deux dés}
\demarche{Pour un minimum, viser P(U supérieur ou égal à k) : les deux dés
doivent dépasser k.}
\begin{questions}
  \item $(U\geqslant k)=(X\geqslant k)\cap(Y\geqslant k)$, et par
    indépendance $P(U\geqslant k)=\Bigl(\dfrac{7-k}{6}\Bigr)^{2}$. D'où
    $P(U=k)=P(U\geqslant k)-P(U\geqslant k+1)=\res{\dfrac{13-2k}{36}}$ :
    la loi charge les petites valeurs.
  \item $\mathrm{E}(U)=\displaystyle\sum_{k=1}^{6}P(U\geqslant k)
    =\frac{36+25+16+9+4+1}{36}=\res{\frac{91}{36}}$ (sommer les queues :
    l'astuce reine pour les variables entières). Et $U+V=X+Y$ donne
    $\mathrm{E}(V)=7-\dfrac{91}{36}=\res{\dfrac{161}{36}}$.
\end{questions}
\end{corrige}
\end{blocLivre}

\begin{blocLivre}{cor-17-12}
\begin{corrige}{12}{Variance d'une somme, covariance}
\demarche{min × max = produit : la covariance de (U,V) se calcule sans leur
loi conjointe.}
\begin{questions}
  \item Par indépendance,
    $\mathrm{E}(XY)=\mathrm{E}(X)\mathrm{E}(Y)=\dfrac{49}{4}$ et
    $\mathrm{V}(X+Y)=2\,\mathrm{V}(X)=\res{\dfrac{35}{6}}$.
  \item $UV=XY$ toujours, donc
    $\mathrm{E}(UV)=\dfrac{49}{4}$ et
    \[
    \operatorname{Cov}(U,V)=\frac{49}{4}-\frac{91}{36}\cdot\frac{161}{36}
    =\res{\frac{1225}{1296}}>0 :
    \]
    un grand minimum va de pair avec un grand maximum — corrélation
    positive, conforme à l'intuition.
\end{questions}
\end{corrige}
\end{blocLivre}

\begin{blocLivre}{cor-17-13}
\begin{corrige}{13}{La binomiale comme somme de Bernoulli}
\demarche{Écrire la binomiale comme somme rend tous ses paramètres
gratuits.}
\begin{questions}
  \item Une issue avec succès aux positions d'une partie $I$ de cardinal
    $k$ a pour probabilité $p^{k}(1-p)^{n-k}$ (indépendance), et il y a
    $\binom nk$ telles parties : $P(S=k)=\binom nk p^{k}(1-p)^{n-k}$.
  \item Linéarité : $\mathrm{E}(S)=np$ ; indépendance (donc
    décorrélation) : $\mathrm{V}(S)=\sum\mathrm{V}(X_i)=\res{np(1-p)}$.
  \item $X+Y$ est une somme de $n+m$ Bernoulli indépendantes (lemme des
    coalitions) : $X+Y\sim\res{\mathcal{B}(n+m,p)}$.
\end{questions}
\end{corrige}
\end{blocLivre}

\begin{blocLivre}{cor-17-14}
\begin{corrige}{14}{Décorrélées mais pas indépendantes}
\demarche{Une covariance nulle ne voit que les moyennes ; l'indépendance
exige bien plus — sauf pour les indicatrices.}
\begin{questions}
  \item $\mathrm{E}(X)=0$ et $\mathrm{E}(XY)=\mathrm{E}(X^{3})=0$ :
    $\operatorname{Cov}(X,Y)=\res{0}$.
  \item $P(X=0,\,Y=1)=0$ alors que
    $P(X=0)P(Y=1)=\dfrac13\cdot\dfrac23\neq0$ : \res{\text{pas
    indépendantes}} — $Y$ est même une fonction de $X$ !
  \item $\mathrm{E}(\ind A\ind B)=P(A\cap B)$ :
    $\operatorname{Cov}(\ind A,\ind B)=P(A\cap B)-P(A)P(B)$, nulle si et
    seulement si $A$ et $B$ sont indépendants : pour les indicatrices,
    décorrélation et indépendance coïncident.
\end{questions}
\end{corrige}
\end{blocLivre}

\begin{blocLivre}{cor-17-15}
\begin{corrige}{15}{Points fixes d'une permutation aléatoire}
\demarche{Découper en indicatrices : l'espérance ignore la dépendance.}
\begin{questions}
  \item $P(\sigma(i)=i)=\dfrac{(n-1)!}{n!}=\dfrac1n$, et par linéarité
    $\mathrm{E}(X)=n\cdot\dfrac1n=\res{1}$ — déjà vu par double comptage
    au chapitre 16, désormais en une ligne.
  \item $X(X-1)$ compte les couples ordonnés $(i,j)$, $i\neq j$, de
    points fixes :
    $\mathrm{E}\bigl(X(X-1)\bigr)=n(n-1)\cdot\dfrac{(n-2)!}{n!}=1$, d'où
    $\mathrm{E}(X^{2})=2$ et $\mathrm{V}(X)=2-1=\res{1}$ : espérance et
    variance valent $1$, quel que soit $n$.
\end{questions}
\end{corrige}
\end{blocLivre}

\begin{blocLivre}{cor-17-16}
\begin{corrige}{16}{Markov et Tchebychev : premiers usages}
\demarche{Des bornes universelles : grossières, mais sans aucune hypothèse
sur la loi.}
\begin{questions}
  \item Markov : $P(X\geqslant10)\leqslant\dfrac{2}{10}=\res{\dfrac15}$.
  \item Tchebychev :
    $P\bigl(|X-3{,}5|\geqslant2\bigr)\leqslant\dfrac{35/12}{4}
    =\dfrac{35}{48}\approx0{,}73$. Exactement : l'événement est
    $X\in\ens{1,6}$, de probabilité $\res{\dfrac13}$. La borne est très
    large — sa force est d'être valable pour \emph{toute} variable de
    même variance.
\end{questions}
\end{corrige}
\end{blocLivre}

\begin{blocLivre}{cor-17-17}
\begin{corrige}{17}{La loi faible des grands nombres}
\demarche{La variance de la moyenne fond en 1/n : Tchebychev fait le reste.}
\begin{questions}
  \item $\mathrm{E}(M_n)=\dfrac12$ et, par indépendance,
    $\mathrm{V}(M_n)=\dfrac{1}{n^{2}}\cdot n\cdot\dfrac14=\res{\dfrac{1}{4n}}$.
  \item Tchebychev :
    $P\bigl(|M_n-\tfrac12|\geqslant\varepsilon\bigr)
    \leqslant\dfrac{1}{4n\varepsilon^{2}}$.
  \item $\dfrac{1}{4n\times0{,}01}\leqslant0{,}05$ donne
    $\res{n\geqslant500}$ lancers. (La vraie probabilité est alors bien
    plus petite encore : la borne est prudente.)
\end{questions}
\end{corrige}
\end{blocLivre}

\begin{blocLivre}{cor-17-18}
\begin{corrige}{18}{Le paradoxe des anniversaires}
\demarche{Passer au complémentaire : tous distincts se calcule en chaîne.}
\begin{questions}
  \item Par les probabilités composées (ou l'équiprobabilité des
    $365^{n}$ listes) : la $k$-ième personne évite les $k-1$ jours déjà
    pris avec probabilité $\dfrac{365-(k-1)}{365}$, d'où le produit.
  \item Pour $n=23$, le produit vaut $\approx0{,}493$ : la probabilité
    d'une coïncidence est $\approx\res{0{,}507}$. Surprenant, mais
    $23$ personnes forment déjà $\binom{23}2=253$ paires — c'est le
    nombre de paires, pas de personnes, qui gouverne. ($n=50$ :
    plus de $96\,\%$.)
\end{questions}
\end{corrige}
\end{blocLivre}

\begin{blocLivre}{cor-17-19}
\begin{corrige}{19}{Dérangements : la loi des points fixes}
\demarche{Choisir les points fixes, déranger le reste, puis laisser n
partir à l'infini.}
\begin{questions}
  \item D'après le chapitre 16,
    $P(X=k)=\dfrac{\binom nk D_{n-k}}{n!}
    =\res{\dfrac{1}{k!}\cdot\dfrac{D_{n-k}}{(n-k)!}}$
    (simplification des factorielles).
  \item $\dfrac{D_m}{m!}\to\mathrm{e}^{-1}$ (chapitre 16), donc à $k$
    fixé :
    $P(X=k)\to\res{\dfrac{\mathrm{e}^{-1}}{k!}}$ — la loi de Poisson de
    paramètre $1$, dont espérance et variance valent $1$ : tout concorde
    avec l'exercice 15.
\end{questions}
\end{corrige}
\end{blocLivre}

\begin{blocLivre}{cor-17-20}
\begin{corrige}{20}{Concentration de la binomiale et sondages}
\demarche{Majorer p(1-p) par 1/4 rend la borne utilisable sans connaître p.}
\begin{questions}
  \item $p(1-p)=\dfrac14-\bigl(p-\tfrac12\bigr)^{2}\leqslant\dfrac14$ ;
    comme $\mathrm{V}(F)=\dfrac{p(1-p)}{n}$, Tchebychev donne
    $P\bigl(|F-p|\geqslant\varepsilon\bigr)
    \leqslant\res{\dfrac{1}{4n\varepsilon^{2}}}$, borne indépendante du
    paramètre inconnu.
  \item $\dfrac{1}{4n\times0{,}0025}\leqslant0{,}05$ exige
    $\res{n\geqslant2000}$ personnes — l'ordre de grandeur des sondages
    réels.
\end{questions}
\end{corrige}
\end{blocLivre}

\begin{blocLivre}{cor-17-21}
\begin{corrige}{21}{Le problème de Monty Hall}
\demarche{La bonne question : quand gagne-t-on en changeant ? Exactement
quand le premier choix était mauvais.}
\begin{questions}
  \item La voiture est placée uniformément ; le choix initial est correct
    avec probabilité $\dfrac13$. En \emph{gardant}, on gagne exactement
    dans ce cas : $P=\dfrac13$. En \emph{changeant}, on gagne exactement
    quand le choix initial était mauvais (l'animateur a alors éliminé
    l'unique autre porte vide, le changement tombe sur la voiture) :
    $P=\res{\dfrac23}$.
  \item L'ouverture n'est pas une information neutre : l'animateur peut
    \emph{toujours} ouvrir une porte vide, quel que soit le choix. Le
    conditionnement par son geste laisse donc
    $P(\text{choix correct})=\frac13$ inchangée — le « 50-50 »
    confondrait avec une ouverture au hasard qui aurait pu révéler la
    voiture.
\end{questions}
\end{corrige}
\end{blocLivre}

\begin{blocLivre}{cor-17-22}
\begin{corrige}{22}{L'espérance du maximum}
\demarche{Le maximum se contrôle par sa fonction de répartition ; sommer
les queues donne l'espérance.}
\begin{questions}
  \item $(M\leqslant k)=\bigcap_i(X_i\leqslant k)$ : par indépendance
    $P(M\leqslant k)=\bigl(\tfrac kN\bigr)^{n}$, et
    \[
    \mathrm{E}(M)=\sum_{k\geqslant1}P(M\geqslant k)
    =\res{\sum_{k=0}^{N-1}\Bigl(1-\Bigl(\frac kN\Bigr)^{n}\Bigr)} .
    \]
  \item $n=2$, $N=6$ :
    $\sum_{k=0}^{5}\bigl(1-\tfrac{k^{2}}{36}\bigr)=6-\tfrac{55}{36}
    =\dfrac{161}{36}$ — l'espérance du maximum de deux dés
    (exercice 11).
  \item $\dfrac{\mathrm{E}(M)}{N}=\dfrac1N\sum_{k=0}^{N-1}
    \Bigl(1-\bigl(\tfrac kN\bigr)^{n}\Bigr)$ est une somme de Riemann :
    elle tend vers
    $\displaystyle\int_{0}^{1}(1-t^{n})\,\mathrm{d}t=\res{\frac{n}{n+1}}$
    (chapitre 15) : avec $n$ tirages, le maximum monte à la fraction
    $\frac{n}{n+1}$ de l'échelle.
\end{questions}
\end{corrige}
\end{blocLivre}

\begin{blocLivre}{cor-17-23}
\begin{corrige}{23}{La marche aléatoire symétrique}
\demarche{Espérance nulle, variance n : la marche vit à l'échelle racine
de n ; le binôme central gouverne les retours.}
\begin{questions}
  \item $\mathrm{E}(S_n)=0$ et, par indépendance
    ($\mathrm{V}(\varepsilon_i)=1$) : $\mathrm{V}(S_n)=\res{n}$.
    Tchebychev : $P\bigl(|S_n|\geqslant a\sqrt n\bigr)\leqslant\dfrac{1}{a^{2}}$ —
    la marche s'éloigne typiquement en $\sqrt n$, pas en $n$.
  \item $S_{2n}=0$ exige exactement $n$ pas $+1$ parmi $2n$ :
    $P(S_{2n}=0)=\dbinom{2n}{n}4^{-n}
    \sim\res{\dfrac{1}{\sqrt{\pi n}}}$ (équivalent du binôme central,
    chapitres 11, 15 et 16) : les retours en $0$ se raréfient comme
    $1/\sqrt n$ — lentement.
\end{questions}
\end{corrige}
\end{blocLivre}

\begin{blocLivre}{cor-17-24}
\begin{corrige}{24}{Weierstrass par Bernstein}
\demarche{Près de x, la continuité uniforme contrôle f ; loin de x, la
concentration binomiale rend les poids négligeables.}
\begin{questions}
  \item $B_n(1)(x)=\sum_k\binom nk x^{k}(1-x)^{n-k}=1$ (binôme), et
    $B_n(t\mapsto t)(x)=\mathrm{E}\bigl(\tfrac{S_n}{n}\bigr)=x$.
  \item Écrivons
    $B_n(f)(x)-f(x)
    =\mathrm{E}\bigl(f(\tfrac{S_n}{n})-f(x)\bigr)$ et séparons selon
    $\bigl|\tfrac{S_n}{n}-x\bigr|<\delta$ ou non. Sur le premier
    événement, $|f(\tfrac{S_n}{n})-f(x)|\leqslant\varepsilon$ (continuité
    uniforme, chapitre 15) ; sur le second, il est majoré par
    $2\max|f|$, et la concentration (exercice 20) donne
    $P\bigl(|\tfrac{S_n}{n}-x|\geqslant\delta\bigr)
    \leqslant\dfrac{1}{4n\delta^{2}}$. En combinant :
    \[
    \bigl|B_n(f)(x)-f(x)\bigr|
    \leqslant\varepsilon+\frac{2\max|f|}{4n\delta^{2}} ,
    \]
    borne \emph{indépendante de $x$}.
  \item Pour $n$ assez grand, le second terme passe sous $\varepsilon$ :
    $\sup_{x}|B_n(f)(x)-f(x)|\leqslant2\varepsilon$. Ainsi
    \res{\text{toute fonction continue sur un segment est limite uniforme
    de polynômes}} — le théorème de Weierstrass, démontré par les
    probabilités : le plus beau pont du programme entre hasard et
    analyse.
\end{questions}
\end{corrige}
\end{blocLivre}

% !TEX root = ../../main.tex
%% Corrigé du défi du chapitre 17 — vérifié par verifications/defis_16_20.py.

\begin{blocLivre}{cor-17-25a}
\begin{corrige}{25}{Le collectionneur de vignettes}
\demarche{Écrire le nombre de vignettes manquantes comme une somme
d'indicatrices : l'espérance est immédiate, la variance aussi, et Markov
puis Tchebychev encadrent le seuil.}
\partieExo{A}{Le nombre moyen de vignettes obtenues}
\begin{questions}
  \item $A_i$ est réalisé par les $(n-1)^{m}$ suites évitant $i$ :
    $P(A_i)=\bigl(1-\frac1n\bigr)^{m}$. Par linéarité,
    $\mathrm{E}(Y_m)=\sum_iP(A_i)=n\bigl(1-\frac1n\bigr)^{m}$ et
    $\res{\mathrm{E}(X_m)=n\Bigl(1-\bigl(1-\tfrac1n\bigr)^{m}\Bigr)}$.
  \item $1-x\leqslant\mathrm{e}^{-x}$ donne
    $\bigl(1-\frac1n\bigr)^{m}\leqslant\mathrm{e}^{-m/n}$. Pour $n=100$ :
    $\mathrm{E}(X_m)\geqslant90\iff0{,}99^{m}\leqslant0{,}1\iff
    m\geqslant\frac{\ln10}{-\ln0{,}99}\approx229{,}1$ : \res{230} achats.
    Les dernières vignettes, elles, coûtent cher\dots
\end{questions}
\partieExo{B}{Le seuil $n\ln n$}
\begin{questions}[resume]
  \item Markov : $P(Y_m\geqslant1)\leqslant\mathrm{E}(Y_m)$. Si
    $m\geqslant n\ln n+cn$ :
    $\mathrm{E}(Y_m)\leqslant n\,\mathrm{e}^{-\ln n-c}=\mathrm{e}^{-c}$,
    d'où $P(\text{complet})\geqslant\res{1-\mathrm{e}^{-c}}$. Pour $n=100$,
    $c=3$ : $761$ achats garantissent l'album complet avec une probabilité
    d'au moins $95\,\%$.
  \item $A_i\cap A_j$ : suites évitant $i$ et $j$, d'où
    $q=P(A_i\cap A_j)=\bigl(1-\frac2n\bigr)^{m}$. Avec $p=\bigl(1-\frac1n\bigr)^{m}$ :
    $\mathrm{E}(Y_m^{2})=\sum_iP(A_i)+\sum_{i\neq j}P(A_i\cap A_j)=np+n(n-1)q$,
    et $\res{\mathrm{V}(Y_m)=np+n(n-1)q-n^{2}p^{2}}$. Comme
    $1-\frac2n\leqslant\bigl(1-\frac1n\bigr)^{2}$, $q\leqslant p^{2}$, d'où
    $\mathrm{V}(Y_m)\leqslant np-np^{2}\leqslant\mathrm{E}(Y_m)$.
  \item $Y_m=0$ impose $|Y_m-\mathrm{E}(Y_m)|\geqslant\mathrm{E}(Y_m)$ :
    Tchebychev donne
    $P(Y_m=0)\leqslant\frac{\mathrm{V}(Y_m)}{\mathrm{E}(Y_m)^{2}}\leqslant\frac{1}{\mathrm{E}(Y_m)}$.
    L'inégalité $\ln(1-x)\geqslant-\frac{x}{1-x}$ (étudier la différence,
    de dérivée $\frac{x}{(1-x)^{2}}\geqslant0$) donne
    $\bigl(1-\frac1n\bigr)^{m}\geqslant\mathrm{e}^{-m/(n-1)}$, donc si
    $m\leqslant(n-1)(\ln n-c)$ : $\mathrm{E}(Y_m)\geqslant\mathrm{e}^{c}$ et
    $P(\text{complet})\leqslant\res{\mathrm{e}^{-c}}$. Juste en dessous de
    $n\ln n$, l'album est presque sûrement incomplet ; juste au-dessus,
    presque sûrement complet : c'est un \emph{phénomène de seuil}.
\end{questions}
\end{corrige}
\end{blocLivre}

\begin{blocLivre}{cor-17-25b}
\begin{corrigeSuite}
\partieExo{C}{Formule exacte pour un petit album}
\begin{questions}[start=6]
  \item Un résultat de $\Omega$ est une application de $\intFF1m$ (les
    achats) dans $\intFF1n$ (les vignettes) ; l'album est complet si et
    seulement si elle est surjective : $n!\,S(m,n)$ cas favorables sur
    $n^{m}$, d'où $\res{\dfrac{n!\,S(m,n)}{n^{m}}}$.
  \item $m=12$ : $\dfrac{720\times1\,323\,652}{6^{12}}
    =\dfrac{953\,029\,440}{2\,176\,782\,336}\approx0{,}438<\frac12$ ;
    $m=13$ : $\dfrac{720\times9\,321\,312}{6^{13}}\approx0{,}514>\frac12$.
    Il faut donc \res{13} lancers. C'est plus que $6\ln6\approx10{,}8$ :
    pour un petit $n$, le terme correctif compte — la médiane vaut environ
    $n\ln n+n\ln\frac{1}{\ln2}\approx n\ln n+0{,}37\,n\approx12{,}95$.
\end{questions}
\end{corrigeSuite}
\end{blocLivre}

